% Lösungen Einheit 4 von 4 – Sinussatz (zusammenbau v0.1)
\einheitenkopf{Einheit 4 von 4 · Sinussatz}

\erg{26}{a) sin, cos, tan \quad b) $\frac{b}{\mathrm{sin}\,51^\circ} = \frac{7}{\mathrm{sin}\,68^\circ}$, $b = 7 \cdot \mathrm{sin}\,51^\circ : \mathrm{sin}\,68^\circ \approx 5{,}9$ cm \quad c) $PQ$ liegt gegenüber $R$, $PR$ gegenüber $Q$: $PR = 8 \cdot \mathrm{sin}\,53^\circ : \mathrm{sin}\,69^\circ \approx 6{,}8$ cm \quad d) $\gamma = 180^\circ - 63^\circ - 48^\circ = 69^\circ$; $c = 8 \cdot \mathrm{sin}\,69^\circ : \mathrm{sin}\,63^\circ \approx 8{,}4$ cm \quad e) $BD$ liegt gegenüber $118^\circ$, $AD$ gegenüber $33^\circ$: $BD = 4{,}6 \cdot \mathrm{sin}\,118^\circ : \mathrm{sin}\,33^\circ \approx 7{,}5$ cm \quad f) $\gamma = 180^\circ - 27^\circ - 118^\circ = 35^\circ$; $c = 7 \cdot \mathrm{sin}\,35^\circ : \mathrm{sin}\,27^\circ \approx 8{,}8$ cm \quad g) $b = 4{,}35 \cdot \mathrm{sin}\,47{,}6^\circ : \mathrm{sin}\,58{,}2^\circ \approx 3{,}8$ m \quad h) $b = 9 \cdot \mathrm{sin}\,44^\circ : \mathrm{sin}\,77^\circ \approx 6{,}42$ cm, also rund $6{,}4$ cm \quad i) $\gamma = 71^\circ$; $a = 10 \cdot \mathrm{sin}\,53^\circ : \mathrm{sin}\,71^\circ \approx 8{,}4$ cm, $b = 10 \cdot \mathrm{sin}\,56^\circ : \mathrm{sin}\,71^\circ \approx 8{,}8$ cm – die Aussage ist falsch. \quad j) $\angle ACB = 90^\circ - 57^\circ = 33^\circ$, $\angle ACD = 91^\circ - 33^\circ = 58^\circ$; $AD = 8 \cdot \mathrm{sin}\,58^\circ : \mathrm{sin}\,76^\circ \approx 7{,}0$ m \quad k) $GS = 3\,200 \cdot \mathrm{sin}\,57^\circ : \mathrm{sin}\,72^\circ \approx 2\,822$ m; Weg $\approx 4\,622$ m $\approx 4{,}6$ km \quad l) Sinussatz: $b = 8 \cdot \mathrm{sin}\,53^\circ : \mathrm{sin}\,67^\circ \approx 6{,}9$ cm. Höhe: $h = 8 \cdot \mathrm{sin}\,53^\circ \approx 6{,}39$ cm, $b = h : \mathrm{sin}\,67^\circ \approx 6{,}9$ cm – gleich \quad m) $\mathrm{sin}\,\beta = 7 \cdot \mathrm{sin}\,68^\circ : 9 \approx 0{,}7211$, $\beta \approx 46{,}1^\circ$ \quad n) $c^2 = 5^2 + 8^2 - 2 \cdot 5 \cdot 8 \cdot \mathrm{cos}\,47^\circ$, $c \approx 5{,}9$ cm \quad o) $\mathrm{cos}\,\gamma = \frac{5^2 + 7^2 - 6^2}{2 \cdot 5 \cdot 7}$, $\gamma \approx 57{,}1^\circ$ \quad p) $h_c = b \cdot \mathrm{sin}\,\alpha$ und $h_c = a \cdot \mathrm{sin}\,\beta$, also $b \cdot \mathrm{sin}\,\alpha = a \cdot \mathrm{sin}\,\beta$; beide Seiten durch $\mathrm{sin}\,\alpha \cdot \mathrm{sin}\,\beta$ teilen. \quad q) $\gamma = 180^\circ - 31^\circ - 113^\circ = 36^\circ$; $BC = 520 \cdot \mathrm{sin}\,31^\circ : \mathrm{sin}\,36^\circ \approx 455{,}6$ m}
\erg{27}{a) Fehler: Seite mit fremdem Winkel gepaart – $b$ gehört zu $\beta$, $a$ zu $\alpha$. Richtig: $b = 8 \cdot \mathrm{sin}\,43^\circ : \mathrm{sin}\,71^\circ \approx 5{,}8$ cm \quad b) Im linken Teildreieck ist $h = b \cdot \mathrm{sin}\,\alpha$, im rechten $h = a \cdot \mathrm{sin}\,\beta$; beides ist dieselbe Strecke, also sind die Terme gleich. \quad c) Winkel bei $C$: $45^\circ$; $AC = 80 \cdot \mathrm{sin}\,63^\circ : \mathrm{sin}\,45^\circ \approx 100{,}8$ m – länger als $100$ m, sie reicht nicht.}
